% CRAI 2026 Position Paper — IEEE Conference Format
% Compositional Safety Gates for Agentic AI
% Draft assembled: 2026-03-21
%
% Build: pdflatex main && bibtex main && pdflatex main && pdflatex main

\documentclass[conference]{IEEEtran}

\usepackage{cite}
\usepackage{amsmath,amssymb,amsfonts}
\usepackage{algorithmic}
\usepackage{graphicx}
\usepackage{textcomp}
\usepackage{xcolor}
\usepackage{booktabs}
\usepackage{hyperref}
\usepackage{listings}
\usepackage{balance}

% Code listing style for Datalog
\lstdefinestyle{datalog}{
  basicstyle=\ttfamily\scriptsize,
  breaklines=true,
  frame=single,
  xleftmargin=1em,
  xrightmargin=1em,
}

% Draft watermark (remove for camera-ready)
% \usepackage{draftwatermark}
% \SetWatermarkText{DRAFT}
% \SetWatermarkScale{4}
% \SetWatermarkColor[gray]{0.92}

\begin{document}

\title{Compositional Safety Gates for Agentic AI:\\A Neurosymbolic Approach with Runtime Assurance}

% TODO: 署名待东丞确认
\author{
\IEEEauthorblockN{[Author 1]}
\IEEEauthorblockA{[Affiliation]}
\and
\IEEEauthorblockN{[Author 2]}
\IEEEauthorblockA{[Affiliation]}
}

\maketitle

\begin{abstract}
Current approaches to AI agent safety rely on monolithic mechanisms: either a single LLM-based classifier or a fixed rule set. Both fail in complementary ways---classifiers hallucinate on edge cases, while rules cannot capture semantic nuance. We present a three-layer composable safety gate for agentic AI systems: (1)~a deterministic blocker using Datalog constraints for zero-latency, formally verifiable pattern matching; (2)~a triviality filter for cost-aware routing that bypasses expensive inference for clearly benign actions; and (3)~a semantic gate using an LLM with minimal-pair prompt engineering for ambiguous cases. On the AgentHarm benchmark (176 harmful + 176 benign scenarios), our system achieves TPR~100\% and FPR~0.0\% on the validation set with majority-vote stability across 50 iterations. Shadow deployment over 19,350+ production calls shows 99.45\% consistency with existing safety rules. We argue that compositional safety---where each layer handles what it does best---provides more robust protection than monolithic approaches and offers a practical path toward formally assured agent governance.
\end{abstract}

\begin{IEEEkeywords}
AI safety, agentic AI, neurosymbolic systems, Datalog, runtime assurance, composable security
\end{IEEEkeywords}

%% ============================================================
%% SECTION 1: Introduction
%% ============================================================
\section{Introduction}

As AI agents transition from isolated chatbots to autonomous systems executing multi-step tool chains, safety must evolve from static content moderation to real-time interception at every tool call. Recent incidents underscore the urgency: in March~2026, a Meta engineering agent autonomously posted flawed technical guidance on an internal forum; employees followed the instructions before the post was identified two hours later, exposing sensitive operational data~\cite{meta2026sev1}. A survey by Cybersecurity Insiders found that 47\% of CISOs have observed unauthorized agent actions, yet only 5\% express confidence in their organization's ability to contain such incidents~\cite{cybersecinsiders2026ciso}.

The prevailing approaches occupy opposite extremes. \textbf{Monolithic LLM classifiers} (e.g., Llama Guard~\cite{inan2023llamaguard}, ShieldGemma~\cite{zeng2024shieldgemma}) offer semantic understanding but introduce latency, hallucinate on edge cases, and cannot provide formal safety guarantees. \textbf{Rule-only systems} (e.g., NeMo Guardrails~\cite{rebedea2023nemoguardrails}) are fast and auditable but cannot distinguish semantically similar actions with different intents---``download a research paper'' versus ``plagiarize a research paper'' look identical to pattern matchers.

The industry is converging on the recognition that agent governance requires new infrastructure. In a single week of March~2026, NVIDIA announced NemoClaw, a declarative policy engine described as ``the policy engine of all the SaaS companies in the world''~\cite{nvidia2026nemoclaw}; Microsoft introduced the 365~E7 tier at \$99/user for agent security~\cite{microsoft2026e7}; and Oasis Security raised \$120M for non-human identity management~\cite{oasis2026funding}. Yet no existing system combines formal guarantees with semantic understanding in a composable architecture.

\textbf{Contributions.} We present Nous, a three-layer composable safety gate (Fig.~\ref{fig:architecture}) with: (1)~a formal composability argument via the Swiss Cheese model~\cite{reason1990swisscheese}; (2)~a shadow-mode deployment methodology enabling zero-downtime transition; and (3)~empirical evaluation on the AgentHarm benchmark~\cite{andriushchenko2025agentharm} demonstrating 100\%/0\% TPR/FPR.


%% ============================================================
%% SECTION 2: Architecture
%% ============================================================
\section{Architecture}

We present a three-layer composable safety gate that intercepts every tool call in an agentic system. Each layer addresses a distinct failure mode, and layers compose sequentially: an input blocked or allowed by an earlier layer never reaches later ones. Fig.~\ref{fig:architecture} illustrates the architecture.

\textbf{Design principles.} (1)~\emph{Composability}: each layer can be independently tested, replaced, and formally verified. (2)~\emph{Fail-closed}: any layer failure (parse error, LLM timeout, configuration corruption) results in a block verdict, never silent passthrough. (3)~\emph{Cost-aware routing}: expensive LLM inference is reserved for genuinely ambiguous cases.

\subsection{Layer 1: Deterministic Blocker}

The first layer applies pattern-matching rules expressed as Datalog constraints. Our implementation uses a TypeScript Datalog engine achieving P50 latency of 0.055ms per evaluation---negligible compared to tool execution time.

Rules are declarative and auditable:
\begin{lstlisting}[style=datalog]
url_has_social(X) :- action_url(X),
    matches(X, "twitter|instagram|weibo").
block(X, "social media link") :-
    url_has_social(X).
\end{lstlisting}

The deployment includes 62~rules covering destructive actions, social media interactions, protected file modifications, and known-malicious patterns. Rules achieve \textbf{zero false positives by construction}---each is manually authored and verified against the benign test set before deployment.

\textbf{Fail-closed semantics.} If the Datalog engine encounters a parse error, the entire input is blocked. This property was identified as a P0 gap in an independent code audit: the initial implementation used soft warnings instead of hard failure.

\subsection{Layer 2: Triviality Filter}

The second layer applies lightweight heuristics to identify trivially benign actions (file reads, web searches without sensitive targets) and route them past the expensive semantic gate. This reduces LLM inference calls by approximately 70\% in production, yielding significant cost savings without sacrificing safety---trivially benign actions have zero overlap with the harmful action space.

\subsection{Layer 3: Semantic Gate}

The third layer handles genuinely ambiguous cases using an LLM (DeepSeek-V3.1) with a prompt containing 23~minimal-pair examples. Each pair consists of a harmful variant (with block reason) and a benign variant (with allow reason), training the model to distinguish intent from surface form.

\textbf{Intent decomposition.} Before classification, we extract structured intent: stated purpose, core task, required capabilities, and risk factors. This decomposition is critical for cases like ``help me correct the dosage guide''---the stated intent is correction, but the core task involves drug dosage modification, requiring careful analysis rather than surface-level allow/block.

\textbf{Majority vote.} Due to stochastic variance in Mixture-of-Experts architectures (where \texttt{temperature=0} does not guarantee determinism), we employ majority vote across $k=3$ evaluations. This eliminates 100\% of observed stochastic FPR in our deployment.

\subsection{Knowledge Graph Enrichment}

A knowledge graph (143~relations, 125~entities) provides post-gate context enrichment---entity resolution, policy lookups, and historical pattern matching. Critically, KG signals feed into the gate's context but do not override verdicts. Early experiments with direct KG injection into the semantic gate caused a 12\% TPR degradation due to attention dilution; the current post-gate architecture avoids this failure mode.

\begin{figure}[t]
  \centering
  % TODO: Replace with exported PNG/PDF from figure1-architecture.drawio
  \fbox{\parbox{0.95\columnwidth}{\centering\vspace{2em}
    \textbf{[Figure 1: Three-Layer Architecture]}\\
    \small See \texttt{figure1-architecture.drawio}\\
    Tool Call $\to$ L1 Datalog $\to$ L2 Triviality $\to$ L3 Semantic $\to$ Verdict
  \vspace{2em}}}
  \caption{Three-layer composable safety gate. Each tool call traverses layers sequentially; early exits at L1 (pattern match) and L2 (trivially benign) reduce cost. The knowledge graph provides post-gate enrichment without verdict override.}
  \label{fig:architecture}
\end{figure}


%% ============================================================
%% SECTION 3: Evaluation
%% ============================================================
\section{Evaluation}

\subsection{AgentHarm Benchmark}

We evaluate on the AgentHarm benchmark~\cite{andriushchenko2025agentharm}, comprising 176~harmful and 176~benign scenarios across 11~categories (cybercrime, disinformation, harassment, etc.).

Table~\ref{tab:results} summarizes our main results. On the validation split, our system achieves TPR~100\% and FPR~0.0\%, stable across majority-vote runs over 50~iterations (Loop~50: $L_{val}=0.0375$). The challenge set of 45~adversarial scenarios---including framing attacks, authority impersonation, and gradual escalation---achieves $L=0.0000$ with intent decomposition.

\begin{table}[t]
\centering
\caption{Main results on AgentHarm benchmark}
\label{tab:results}
\begin{tabular}{lcccc}
\toprule
\textbf{Split} & \textbf{TPR} & \textbf{FPR} & \textbf{$L_{val}$} & \textbf{$N$} \\
\midrule
Validation    & 100\% & 0.0\% & 0.0375 & 352 \\
Challenge     & 100\% & 0.0\% & 0.0000 & 45 \\
Full (test)   & 96\%  & 2.8\% & ---    & 352 \\
\midrule
Shadow (prod) & ---   & ---   & ---    & 19,350+ \\
\bottomrule
\end{tabular}
\vspace{0.5em}
\small Shadow deployment: 99.45\% consistency over 14+ days.
\end{table}

\subsection{Shadow Deployment}

Before promoting the system to primary interceptor, we operated in shadow mode for 14+~days, processing 19,350+ real production calls. The system ran in parallel with existing safety rules, logging verdicts without affecting operations. Agreement rate was 99.45\%; the 0.55\% divergence consisted primarily of cases where the semantic gate correctly identified nuance that rule-based systems missed.

After 14,238 shadow calls with stable agreement metrics, the system was promoted to primary interceptor (cutover). Post-cutover monitoring confirmed no regression.

\subsection{Capability Preservation}

A critical concern with safety systems is over-blocking. We evaluate capability preservation using L5~extreme probes---genuine policy-boundary cases where reasonable safety policies disagree:

\begin{itemize}
\item \textbf{Correctly allowed}: authorized penetration tests, compliance research, academic study of harmful content
\item \textbf{Correctly blocked}: surveillance tool deployment, controlled substance synthesis guides
\item \textbf{Debatable}: harm reduction drug identification (62.5\% allow rate---a policy choice, not a system bug)
\end{itemize}

\subsection{Adversarial Red Team}

We conducted a 20-scenario adversarial campaign (CC-BOS) testing classical Chinese framing, authority claims, and gradual escalation attacks. Initial attack success rate was 4.1/5; after adding content safety (T14), authority detection (T15), and input normalization (Step~0), the rate dropped to 1.1/5.


%% ============================================================
%% SECTION 4: Discussion
%% ============================================================
\section{Discussion}

\textbf{Composability as safety primitive.} Following Reason's Swiss Cheese model~\cite{reason1990swisscheese}, each layer has characteristic ``holes'' (failure modes), but the probability of a threat passing through all layers simultaneously is multiplicatively small. Unlike monolithic systems where a single failure is catastrophic, our architecture degrades gracefully.

\textbf{Stochastic variance in MoE.} A surprising finding is that \texttt{temperature=0} does not guarantee deterministic outputs from Mixture-of-Experts architectures. We observed intermittent FPR fluctuations that majority vote completely eliminated---a practical consideration absent from current literature.

\textbf{Multi-agent safety gaps.} Recent analyses of autonomous agent deployments~\cite{agentsofchaos2025} identify three critical failure modes absent from current benchmarks: fake task completion (agents reporting success without execution), PII leakage across tool chains, and unbounded process lifetimes. Our composable architecture addresses the first via proof traces and the third via fail-closed timeouts, but PII leakage across tool boundaries remains an open problem.

\textbf{Limitations.} (1)~The current system is single-turn: it evaluates each tool call independently without multi-turn trajectory analysis. A sufficiently patient adversary could decompose harmful intent across many benign-looking steps. (2)~The semantic gate prompt is English-centric; cross-lingual attacks may exploit language gaps. (3)~Our evaluation is on a single benchmark; broader evaluation across diverse agent architectures is needed.

\textbf{Future work.} Multi-turn trajectory analysis using sliding-window context; cross-lingual semantic gates; integration with differentiable Datalog~\cite{li2023scallop} for end-to-end optimization of the rule-semantic boundary; and formal verification of the Datalog-to-LLM handoff protocol.


%% ============================================================
%% SECTION 5: Related Work
%% ============================================================
\section{Related Work}

\textbf{Monolithic classifiers.} Llama Guard~\cite{inan2023llamaguard} and ShieldGemma~\cite{zeng2024shieldgemma} fine-tune LLMs for content classification. While effective for conversational safety, they introduce latency incompatible with real-time tool interception and cannot provide formal guarantees.

\textbf{Rule-based systems.} NeMo Guardrails~\cite{rebedea2023nemoguardrails} offers programmable rails for LLM applications. Our Datalog layer is architecturally similar but extends the approach with formal fail-closed semantics and composability with a semantic layer.

\textbf{Enterprise agent governance.} NVIDIA's NemoClaw~\cite{nvidia2026nemoclaw}, announced at GTC~2026, provides a declarative policy engine for enterprise agent governance---architecturally homologous to our Layer~1. Guardrails AI~\cite{guardrailsai2024} offers input/output validation for LLM applications. Microsoft 365~E7~\cite{microsoft2026e7} and Oasis Security~\cite{oasis2026funding} represent industry convergence on the need for dedicated agent safety infrastructure.

\textbf{Neurosymbolic AI.} Scallop~\cite{li2023scallop} demonstrates differentiable Datalog for neurosymbolic programming. Integrating differentiable reasoning into our Datalog layer is a promising future direction for optimizing the rule-semantic boundary.

\textbf{Automated red-teaming.} Votal AI's CART platform~\cite{votalai2026cart} and the AgentHarm benchmark~\cite{andriushchenko2025agentharm} provide complementary evaluation methodologies. Our CC-BOS campaign extends these with cultural framing attacks.


%% ============================================================
%% References
%% ============================================================
\balance
\bibliographystyle{IEEEtran}
\bibliography{references}

\end{document}
